\section{Introduction}

We study the polynomial family $\{Z_m(y)\}_{m\ge 0}$ defined by the rational generating function
\[
Z(t,y)=\sum_{m\ge 0} Z_m(y)t^m
=
\frac{1+yt+(y^2-y-1)t^2+(y^3-2y)t^3}
{1-t-(2y+1)t^2+t^3+y(y+1)t^4}.
\]
Equivalently, $\{Z_m(y)\}$ is generated by an order-$4$ companion transfer matrix with characteristic polynomial
\[
\Pi(\lambda,y)=\lambda^4-\lambda^3-(2y+1)\lambda^2+\lambda+y(y+1).
\]
The paper treats this quartic as a fixed algebraic-spectral object and asks why it is the canonical member to analyze inside its four-parameter family.
Our point of view is purely spectral: we track the elliptic normalization of the projective quartic, the exact relation between the dominant-modulus envelope and the ramification locus, and the effect of the dominant square-root branch point on the zeros of \(Z_m\) nearest $y=0$.
Section~\ref{sec:genus-one} proves that \(\Pi(\lambda,y)\) belongs to the balanced family
\[
P_{a,c,d,g}(\lambda,y)=\lambda^4+a\lambda^3+(c-2y)\lambda^2+d\lambda+y(y-c)+g,
\]
for which the substitution \(\xi=y-\lambda^2\) removes the quartic obstruction.
We upgrade Proposition~\ref{prop:balanced-family} by proving a balanced-family classification theorem (Theorem~\ref{thm:balanced-classification}): \((a,c,d,g)=(-1,-1,1,0)\) is the unique integer-balanced member (up to the natural scaling symmetry) with smooth genus one, three finite real branch values, and a dominant square-root branch at $y=0$.
This canonical property supplies the promised motivation requested by the referee and justifies our detailed focus on this particular quartic.
Specializing those balanced-family statements gives the spectral polynomial
\[
Y^2=4\lambda^3-4\lambda+1,
\]
the normalization of the projective quartic.
The four eigenvalue branches thus form a genus-one degree-$4$ cover of the \(y\)-line.
For quartic spectral curves outside this balanced locus the same level of analytic control appears out of reach.

The present quartic is the smooth genus-one specialization \((a,c,d,g)=(-1,-1,1,0)\), so its projective closure is normalized by the elliptic curve
\[
Y^2=4\lambda^3-4\lambda+1,
\]
and this turns the problem into a genus-one covering of the $y$-plane.
As a result, one can compute the finite branch values, classify the real branch points, and write the edge zeros near $y=0$ in closed form.

\subsection{Relation to existing literature}

The algebraic ingredients used in this paper are classical.
The Yang--Lee papers and Fisher's edge-singularity paper provide historical motivation for zero distributions in parameter-dependent polynomial families \cite{YangLee1952I,LeeYang1952II,Fisher1978}, while the modern BKW formulations of Brown and Otto \cite{BrownOtto2020} keep track of how the limiting equimodular sets evolve.
Companion matrices and cyclicity are classical \cite{HornJohnson2012}; rational generating functions, linear recurrences, and finite Hankel rank belong to the general theory of rational sequences \cite{FlajoletSedgewick2009,Stanley2012,Partington1989}; and the reduction to Weierstrass form, the Riemann--Hurwitz count, and the local Puiseux expansions are standard algebraic-curve methods \cite{Walker1950,Silverman2009,Miranda1995,BrieskornKnorrer1986}.
The four-term-recurrence literature most relevant to this paper includes de Bruin's early analysis of zero locations \cite{deBruin1984}, Tran's discriminant approach \cite{Tran2014,Tran2015}, the Forg\'acs--Tran program \cite{ForgacsTran2016,ForgacsTran2017,ForgacsTran2020}, and the more recent works of Tran--Zumba and Botta--Fornazieri that treat linear-coefficient families \cite{TranZumba2018,TranZumba2021,BottaFornazieri2023}.
Wider Toeplitz and banded-matrix perspectives are represented by Ndikubwayo \cite{Ndikubwayo2022} and by the geometric limit analysis of Wang and Zhang for polynomial sequences of type $(1,1)$ \cite{WangZhang2021}.
Compared with those works, the contribution here is not new machinery, but the explicit working-out of a specific quartic case $\Pi(\lambda,y)$ together with an accompanying balanced-family classification: the projective normalization, the factorization of the ramification quintic, the real-branch classification, the real-spectral phase diagram, the dominant-branch identification on $y>0$, the endpoint zero quantization at $y=0$, the degree law, and (in Section~\ref{sec:genus-one}) two precise balanced-family criteria (smoothness plus real projection) that extend beyond the single specialization.

\subsubsection*{What is new relative to the literature}

The closest neighboring viewpoints are Biggs's equimodular-curve picture, Tran's discriminant analysis of rational generating functions, and the four-term-recurrence papers of Forg\'acs--Tran, Tran--Zumba, and Adams.
The point of the following comparison is to calibrate the manuscript as a detailed worked example centered on one quartic, not as a new general theory of short recurrences.

\textbf{Biggs~\cite{Biggs2002}} emphasizes equimodular curves as the relevant boundaries for zero accumulation in recursively defined polynomial families.
That viewpoint is close in spirit to the dominant-modulus discussion used here, but Biggs does not address the algebraic normalization of the quartic spectral curve or the exact real-branch classification for the present family.

\textbf{Brown--Otto~\cite{BrownOtto2020}} describe modern BKW-type formulations that emphasize structural comparisons between spectral polynomials.
Their framework motivates our notation for dominant branches, but they do not construct explicit elliptic normalizations or the balanced-family classification used below.

\textbf{Tran~\cite{Tran2014,Tran2015}} analyzes broader classes of rational-generating-function families via discriminant conditions and identifies real algebraic support sets.
The present quartic falls within Tran's general framework, but the balanced-family reduction and the elliptic normalization are not produced by Tran's discriminant approach, which does not descend to genus-one covers or identify the dominant branch explicitly.

\textbf{de Bruin~\cite{deBruin1984}} treats four-term recurrences via orthogonal-polynomial methods and obtains zero-location results by Sturm theory.
The present work differs in that it relies on the balanced-family elliptic reduction to isolate the dominant branch and to write the endpoint zeros explicitly.

\textbf{Forg\'acs--Tran~\cite{ForgacsTran2016,ForgacsTran2017,ForgacsTran2020}} develop eventual hyperbolicity criteria and zero-distribution results for polynomials with rational generating functions and hyperbolic-type denominators.
The present quartic denominator $D(t,y)$ does not fall into the hyperbolic class treated in those papers (it is not eventually hyperbolic for all real $y$), and those papers do not address the dominant-modulus identification or the endpoint zero asymptotics for the nonhyperbolic regime.

\textbf{Tran--Zumba~\cite{TranZumba2018,TranZumba2021}} study zeros of polynomials with four-term recurrences, including the case of linear coefficients.
The present quartic family has constant coefficients (in $m$) but a one-parameter $y$-deformation, so it is a different kind of four-term problem.
Those papers describe real-zero conditions for specific coefficient patterns but do not produce an elliptic normalization or a branch-value phase diagram.

\textbf{Adams~\cite{Adams2020}} treats hyperbolic polynomials with four-term recurrence and linear coefficients (Calcolo 2020).
That setting is different from ours: Adams considers hyperbolicity of individual polynomials in their coefficient sequence, whereas the present paper concerns the $y$-parametric spectral polynomial $\Pi(\lambda,y)$ and its elliptic reduction.

\textbf{Ndikubwayo~\cite{Ndikubwayo2022}} extends the recurrence framework to five-term recurrences, banded Toeplitz matrices, and reality of zeros.
The present quartic is a four-term system and its spectral analysis rests on different structure (genus-one reduction rather than Toeplitz machinery).

\textbf{Botta--Fornazieri~\cite{BottaFornazieri2023}} study four-term recurrences with linear coefficients and obtain zero-location information via eigenvalue methods.
Their classification is parameteric in the recurrence index rather than in an external spectral parameter, so it does not produce the genus-one covering or the dominant-branch inequalities supplied here.

\textbf{Wang--Zhang~\cite{WangZhang2021}} analyze geometric limits of zeros for polynomial sequences of type $(1,1)$.
Their work shares the geometric viewpoint but focuses on Toeplitz determinants instead of balanced quartics; no elliptic normalization is available in their setting.

In each case, what the present paper adds for the specific quartic $\Pi(\lambda,y)$ that none of the above provides is: an explicit elliptic normalization of the projective quartic curve, an exact real branch-value phase diagram from the elliptic real locus, a direct proof of global dominance on the positive axis via explicit modulus inequalities, and the quantitative endpoint zero asymptotics from the local cosine law.
The general balanced-family real-projection criterion (Theorem~\ref{thm:balanced-real-projection}) is the one statement that goes beyond the single specialization.
The correct calibration of the contribution is: an exact elliptic normalization and dominant-branch analysis for one quartic lying outside the previously studied families, together with a balanced-family structural observation that may extend further.

This scope distinction also calibrates the novelty of the individual results.
The minimal-order/Hankel-rank statement of Section~\ref{sec:hankel} is a family-specific specialization of standard reduced-denominator reasoning.
Section~\ref{sec:genus-one} packages the balanced-family reduction, the explicit elliptic normalization, and the real projection criterion for the balanced family.
Sections~\ref{sec:ep-classification} and \ref{sec:physical-sheet} use that normalization to determine the real branch-value picture, to prove dominance on the positive axis by direct modulus inequalities, and to extract the endpoint zero asymptotics.
Section~\ref{sec:support-law} closes with an explicit recurrence computation for degree growth.
The paper should therefore be read as a self-contained quartic case study together with the balanced-family structural observations of Proposition~\ref{prop:balanced-family} and Theorem~\ref{thm:balanced-real-projection}; the remaining theorems are specific to the specialization \((a,c,d,g)=(-1,-1,1,0)\).

The scope of the present article is therefore precise.
For real $y>0$ we write
\[
r(y):=\max_{1\le j\le 4}\abs{\lambda_j(y)},
\qquad
f_{\mathrm{dom}}(y):=\log r(y).
\]
The dominant branch over $y>0$ is the connected component of the smooth spectral curve on which $\abs{\lambda}=r(y)$.
We do not derive the quartic family from a separate microscopic model within this paper.
Rather, we treat the induced companion family as a fixed algebraic-spectral object and ask what can be proved explicitly about its spectral curve, its real branch points, and the zeros of the associated polynomial sequence.

\subsection{Summary of results}

For this quartic family we obtain the following.

\begin{enumerate}
\item \textbf{Balanced-family reduction, canonical classification, and real projection criterion} (Section~\ref{sec:genus-one}).
We identify a four-parameter balanced family \(P_{a,c,d,g}\) for which the substitution \(\xi=y-\lambda^2\) converts the quartic curve to a Weierstrass model, sharpen the analysis by proving Theorem~\ref{thm:balanced-classification} (the canonical classification requested by the referee), and give a general real projection criterion (Theorem~\ref{thm:balanced-real-projection}) determining, for any balanced family member with real parameters, whether the real $y$-image consists of one or two disjoint intervals.
The present transfer polynomial is the specialization \((a,c,d,g)=(-1,-1,1,0)\), with Weierstrass invariants \(g_2=4\), \(g_3=-1\), and \(\Delta=37\).
The resulting elliptic curve is the normalization of the projective quartic, and its finite branch values are \(y=0\), \(y=1\), and the three roots of \(256y^3+411y^2+165y+32\).

\item \textbf{Real branch-point classification and real-root counts} (Section~\ref{sec:ep-classification}).
Among the three real branch values, $y=0$ is the only one whose outgoing pair contains the dominant root for $y>0$ small; $y=1$ and the negative branch value $y_{\mathrm{sub}}$ are subdominant.
The real locus of the elliptic curve projects exactly to $[y_{\mathrm{sub}},1]\cup[0,\infty)$, which yields the full count of real roots of $\Pi(\cdot,y)$ for every real $y$.

\item \textbf{Square-root branching and endpoint quantization} (Section~\ref{sec:physical-sheet}).
We state a generic branching lemma for analytic double roots of a spectral polynomial.
For the present family we verify the amplitude symmetry and complex-uniform remainder bounds needed for the local cosine law, obtaining
\[
y_{m,k}
=
-\frac{\pi^2(2k-1)^2}{2m^2}
-\frac{3\pi^2(2k-1)^2}{4m^3}
+O_k(m^{-4})
\]
for the negative zeros nearest $y=0$.

\item \textbf{Large-$y$ splitting and degree growth} (Sections~\ref{sec:physical-sheet} and \ref{sec:support-law}).
On the positive large-$y$ branch, the dominant root has a quarter-power splitting governed by a double tangency at projective infinity.
On the polynomial side, $\deg_y Z_m=\lfloor(m+3)/2\rfloor$, and the top-degree coefficients are explicit and subexponential in $m$.

\item \textbf{A rigid boundary value at $y=-1$} (Section~\ref{sec:boundary-zeros}).
For every $m\ge 2$, the polynomial $Z_m(y)$ has a zero of multiplicity at least two at $y=-1$.
This gives a uniform algebraic boundary point for the entire sequence.
\end{enumerate}

Sections~\ref{sec:model}--\ref{sec:hankel} define the family and establish the generic order-$4$ minimality of the recurrence.
Sections~\ref{sec:genus-one}--\ref{sec:boundary-zeros} contain the explicit spectral and algebraic results for this quartic family.
Section~\ref{sec:verification} records reproducible numerical checks of the real-axis modulus ordering and the first edge zeros, and Section~\ref{sec:discussion-new} concludes with the scope and limitations of the quartic analysis.

\paragraph{Novelty calibration by result.}
Relative to the works cited above, Theorem~\ref{thm:minimal-order} and the Hankel-rank statement of Section~\ref{sec:hankel} are standard consequences of rational-generating-function theory \cite{FlajoletSedgewick2009,Stanley2012}.
Proposition~\ref{prop:balanced-family} and Theorem~\ref{thm:balanced-classification} provide the new family-level structural contributions: the elliptic normalization and the canonicality of the present quartic inside the balanced family.
The explicit elliptic normalization (Theorem~\ref{thm:elliptic-reduction}), the real-branch classification (Theorem~\ref{thm:ep-classification}), the real-spectral phase diagram (Theorem~\ref{thm:real-phase-diagram}), the dominant-branch identification (Proposition~\ref{prop:global-dominance}), and the endpoint zero asymptotics (Theorem~\ref{thm:vacuum-zero-quantization}) are specialization-specific statements not covered by \cite{Tran2014,Tran2015,deBruin1984,BottaFornazieri2023}.
The degree law (Theorem~\ref{thm:support-law}) and the $y=-1$ rigidity (Proposition~\ref{prop:double-zero}) arise from explicit recurrence manipulations particular to this transfer matrix.
